% Options for packages loaded elsewhere
\PassOptionsToPackage{unicode}{hyperref}
\PassOptionsToPackage{hyphens}{url}
\documentclass[
]{article}
\usepackage{xcolor}
\usepackage{amsmath,amssymb}
\setcounter{secnumdepth}{-\maxdimen} % remove section numbering
\usepackage{iftex}
\ifPDFTeX
  \usepackage[T1]{fontenc}
  \usepackage[utf8]{inputenc}
  \usepackage{textcomp} % provide euro and other symbols
\else % if luatex or xetex
  \usepackage{unicode-math} % this also loads fontspec
  \defaultfontfeatures{Scale=MatchLowercase}
  \defaultfontfeatures[\rmfamily]{Ligatures=TeX,Scale=1}
\fi
\usepackage{lmodern}
\ifPDFTeX\else
  % xetex/luatex font selection
\fi
% Use upquote if available, for straight quotes in verbatim environments
\IfFileExists{upquote.sty}{\usepackage{upquote}}{}
\IfFileExists{microtype.sty}{% use microtype if available
  \usepackage[]{microtype}
  \UseMicrotypeSet[protrusion]{basicmath} % disable protrusion for tt fonts
}{}
\makeatletter
\@ifundefined{KOMAClassName}{% if non-KOMA class
  \IfFileExists{parskip.sty}{%
    \usepackage{parskip}
  }{% else
    \setlength{\parindent}{0pt}
    \setlength{\parskip}{6pt plus 2pt minus 1pt}}
}{% if KOMA class
  \KOMAoptions{parskip=half}}
\makeatother
\usepackage{graphicx}
\makeatletter
\newsavebox\pandoc@box
\newcommand*\pandocbounded[1]{% scales image to fit in text height/width
  \sbox\pandoc@box{#1}%
  \Gscale@div\@tempa{\textheight}{\dimexpr\ht\pandoc@box+\dp\pandoc@box\relax}%
  \Gscale@div\@tempb{\linewidth}{\wd\pandoc@box}%
  \ifdim\@tempb\p@<\@tempa\p@\let\@tempa\@tempb\fi% select the smaller of both
  \ifdim\@tempa\p@<\p@\scalebox{\@tempa}{\usebox\pandoc@box}%
  \else\usebox{\pandoc@box}%
  \fi%
}
% Set default figure placement to htbp
\def\fps@figure{htbp}
\makeatother
% definitions for citeproc citations
\NewDocumentCommand\citeproctext{}{}
\NewDocumentCommand\citeproc{mm}{%
  \begingroup\def\citeproctext{#2}\cite{#1}\endgroup}
\makeatletter
 % allow citations to break across lines
 \let\@cite@ofmt\@firstofone
 % avoid brackets around text for \cite:
 \def\@biblabel#1{}
 \def\@cite#1#2{{#1\if@tempswa , #2\fi}}
\makeatother
\newlength{\cslhangindent}
\setlength{\cslhangindent}{1.5em}
\newlength{\csllabelwidth}
\setlength{\csllabelwidth}{3em}
\newenvironment{CSLReferences}[2] % #1 hanging-indent, #2 entry-spacing
 {\begin{list}{}{%
  \setlength{\itemindent}{0pt}
  \setlength{\leftmargin}{0pt}
  \setlength{\parsep}{0pt}
  % turn on hanging indent if param 1 is 1
  \ifodd #1
   \setlength{\leftmargin}{\cslhangindent}
   \setlength{\itemindent}{-1\cslhangindent}
  \fi
  % set entry spacing
  \setlength{\itemsep}{#2\baselineskip}}}
 {\end{list}}
\usepackage{calc}
\newcommand{\CSLBlock}[1]{\hfill\break\parbox[t]{\linewidth}{\strut\ignorespaces#1\strut}}
\newcommand{\CSLLeftMargin}[1]{\parbox[t]{\csllabelwidth}{\strut#1\strut}}
\newcommand{\CSLRightInline}[1]{\parbox[t]{\linewidth - \csllabelwidth}{\strut#1\strut}}
\newcommand{\CSLIndent}[1]{\hspace{\cslhangindent}#1}
\setlength{\emergencystretch}{3em} % prevent overfull lines
\providecommand{\tightlist}{%
  \setlength{\itemsep}{0pt}\setlength{\parskip}{0pt}}
\usepackage{bookmark}
\IfFileExists{xurl.sty}{\usepackage{xurl}}{} % add URL line breaks if available
\urlstyle{same}
\hypersetup{
  pdftitle={World Tubes in Gauged Camera Space: Frame-Amortized Dynamic Gaussian Rendering},
  pdfauthor={Anonymous},
  hidelinks,
  pdfcreator={LaTeX via pandoc}}

\title{World Tubes in Gauged Camera Space: Frame-Amortized Dynamic
Gaussian Rendering}
\author{Anonymous}
\date{}

\begin{document}
\maketitle

\subsection{Abstract}\label{abstract}

Dynamic Gaussian renderers repeatedly evaluate, project, bin, order, and
differentiate scene primitives for each requested time. This is
efficient for an isolated frame but duplicates world-side work for video
playback, finite exposure, rolling shutter, and repeated supervision
along a known camera path. We introduce \textbf{World Tubes}, a
camera-program compiler that maps dynamic Gaussian primitives into
reusable traces over sensor space and time. The construction pulls a
spacetime primitive through the camera-ray bundle and pushes it forward
along the ray-depth fiber. A local Schur complement gives an exact
Gaussian footprint and conditional-depth packet in an affine chart;
homogeneous projective traces, gauge Jacobians, and event-certified
domains extend the construction to bounded moving-camera segments
without silently crossing projective poles or depth-order events. A
fixed-topology adjoint maps image residuals through the compiled
evaluator to the original world parameters. On certified fixtures,
ordinary- and log-depth gauges agree to \(3.50\times10^{-13}\) in value
and \(2.33\times10^{-12}\) in gradient, while omitting the fiber
Jacobian produces errors above \(0.59\). Visibility stratification
reduces a deliberate order-crossing error from \(0.187\) to zero, and
fixed trace state reaches a \(0.03125\) trace-count ratio relative to
per-frame replay at \(F=128\). These results establish bounded
correctness and structural reuse while retaining the unavoidable cost of
producing the requested sensor samples.

\subsection{1. Introduction}\label{introduction}

Three-dimensional Gaussian splatting replaces neural ray marching with
visibility-aware anisotropic splats and makes high-quality novel-view
rendering interactive (\citeproc{ref-kerbl2023}{Kerbl et al. 2023}).
Dynamic extensions add deformation fields, persistent trajectories,
temporal features, or native spacetime covariance
(\citeproc{ref-wu2024_4dgs}{Wu et al. 2024};
\citeproc{ref-yang2024_deformable}{Ziyi Yang et al. 2024};
\citeproc{ref-luiten2024}{Luiten et al. 2024};
\citeproc{ref-li2024_spacetime}{Li et al. 2024};
\citeproc{ref-yang2024_native}{Zeyu Yang et al. 2024}). These
representations differ in how they model the world, but their renderers
generally retain the same computational unit: for every target time,
evaluate the dynamic primitive state, project it, determine support and
tile membership, establish a visibility order, composite, and
backpropagate through that render.

That frame-by-frame unit repeats work when many samples follow a known
camera program. A video renderer requests neighboring times. A
finite-exposure image uses multiple shutter samples. Rolling shutter
couples row position to capture time. Training repeatedly visits the
same path or a low-dimensional family of paths. The image values still
require \(O(FHW)\) output samples, but the world-side projection,
conservative support, binning, visibility metadata, and their backward
tape need not be rebuilt independently at every time.

World Tubes changes the unit of compilation from a frame to a camera
program. Let \(B=\Omega\times\mathcal T\) be the sensor-time base and
let \(y=(u,v,\tau)\). A dynamic world primitive is compiled into trace
fields

\[
  \bigl(\alpha_i(y),\,c_i(y),\,\widehat z_i(y),\,\nu_{z,i}(y)\bigr),
\]

together with support bounds, visibility certificates, fallback
metadata, and a differentiable map back to the original world
parameters. Frames are slices of this atlas. Exposure and
rolling-shutter images are integrals or row-dependent samples of the
same object.

The key difficulty is not drawing a smooth curve through projected
centers. Perspective division can cross a pole; projected support can
enter a tile; and two overlapping translucent primitives can exchange
depth order. A trace may therefore remain algebraically smooth while a
renderer compiled under one support or order topology becomes incorrect.
We retain the camera-gauge mathematics that detects these events.
Homogeneous traces are divided only on certified projective domains, the
ray-fiber measure includes its coordinate Jacobian, and
support-overlapping traces are partitioned where their order certificate
fails. Rejected domains are split or routed to an explicit fallback
rather than extrapolated.

Our contributions are:

\begin{enumerate}
\def\labelenumi{\arabic{enumi}.}
\tightlist
\item
  A camera-ray-bundle formulation in which the observable trace is the
  fiber pushforward \(\pi_*\Gamma^*w_i\), with an exact local Gaussian
  marginal and conditional-depth packet obtained by a Schur complement.
\item
  A projective trace and visibility atlas whose denominator,
  approximation, support, and order certificates define bounded reuse
  domains under moving cameras.
\item
  A fixed-topology evaluator and compiled adjoint that reuse trace and
  tile-time state while returning gradients to the same learned world.
\item
  A causal evaluation contract that compares per-frame replay and
  compilation from one frozen representation, separating world-side
  reuse from the unavoidable cost of materializing output pixels.
\end{enumerate}

\begin{figure}
\centering
\includegraphics[width=1\linewidth,height=\textheight,keepaspectratio]{/Users/nicholasbardy/git/gsplats_browser/dynaworld/tmp/pdfs/world_tubes_system_overview.png}
\caption{World Tubes compiles projective traces, certified tile--time
cells, and their adjoint once for a bounded camera program. Per-frame
replay repeats the same world-side lowering. Both routes still shade and
materialize every requested sensor
sample.}\label{fig:world-tubes-system-overview}
\end{figure}

\subsection{2. Related work}\label{related-work}

\textbf{Gaussian and dynamic splatting.} Classical surface and volume
splatting derive filtered screen footprints in a locally affine
ray-space approximation (\citeproc{ref-zwicker2001_surface}{Zwicker et
al. 2001b}, \citeproc{ref-zwicker2001_ewa_volume}{2001a}). Modern 3D
Gaussian splatting learns anisotropic primitives with differentiable
visibility-aware rasterization (\citeproc{ref-kerbl2023}{Kerbl et al.
2023}), and Mip-Splatting makes the sampling boundary explicit
(\citeproc{ref-yu2024_mipsplatting}{Yu et al. 2024}). Dynamic methods
extend the representation through deformation, motion, temporal opacity,
or four-dimensional covariance (\citeproc{ref-wu2024_4dgs}{Wu et al.
2024}; \citeproc{ref-yang2024_deformable}{Ziyi Yang et al. 2024};
\citeproc{ref-luiten2024}{Luiten et al. 2024};
\citeproc{ref-li2024_spacetime}{Li et al. 2024};
\citeproc{ref-yang2024_native}{Zeyu Yang et al. 2024}). World Tubes is
complementary: it compiles repeated rendering work induced by a camera
program rather than proposing another dynamic scene family.

\textbf{Nonlinear cameras and temporal sensors.} Gaussian Splatting on
the Move models blur and rolling shutter, while 3DGUT replaces a
first-order projection with an unscented transform for nonlinear cameras
and secondary rays (\citeproc{ref-seiskari2024}{Seiskari et al. 2025};
\citeproc{ref-wu2025_3dgut}{Wu et al. 2025}). These methods improve
projection at a requested sample. Our projective domains instead
determine when a collection of samples can share trace, support, and
visibility state.

\textbf{Visibility and compositing.} Order-independent transparency and
sort-free Gaussian renderers change the transfer law to reduce sorting
(\citeproc{ref-mcguire2013_weighted_oit}{McGuire and Bavoil 2013};
\citeproc{ref-hou2025_sortfree}{Hou et al. 2025};
\citeproc{ref-koo2026_blending}{Koo et al. 2026}). StopThePop improves
view consistency through finer per-ray sorting
(\citeproc{ref-radl2024_stopthepop}{Radl et al. 2024}). World Tubes
preserves the frozen baseline's ordered front-to-back law on certified
cells; it splits or falls back when that noncommutative order cannot be
reused.

\subsection{3. World Tubes}\label{world-tubes}

\subsubsection{3.1 Camera-ray pushforward and
gauges}\label{camera-ray-pushforward-and-gauges}

Let \(\pi:E_\Gamma\rightarrow B\) be the ray bundle of a known camera
program. The fiber \(F_y=\pi^{-1}(y)\) contains the physical ray-depth
points associated with sensor-time sample \(y\). The camera program maps
bundle points into world spacetime,

\[
  \Gamma:E_\Gamma\rightarrow\mathbb R^3\times\mathbb R.
\]

For a world primitive \(w_i\), its sensor-time trace is

\[
  T_i=\pi_*\Gamma^*w_i.
  \tag{1}
\]

Choose a local depth coordinate \(z_a\) on a regular bundle chart
\(C_a\). For a density-like primitive \(\rho_i\), Eq. (1) becomes

\[
  T_i(y)=\int_{D_a}
  \rho_i\!\left(\Gamma_a(y,z_a)\right)
  J_a(y,z_a)\,\mathrm dz_a,
  \tag{2}
\]

where \(J_a\) converts coordinate length to the physical fiber measure.
If \(z_b=h(y,z_a)\) is an orientation-preserving reparameterization of
the same fiber segment, the one-dimensional change-of-variables formula
makes Eq. (2) invariant. A chart is invalid when the coordinate map
becomes singular or the two charts cover different clipped fiber
segments; invariance does not justify crossing a pole, support event, or
visibility event without a transition.

This distinction is operational. Ordinary depth, log depth, inverse
depth, or an orbit-adapted coordinate may describe the same fiber, but
only with the corresponding \(J_a\). Our certified gauge fixture obtains
value and gradient agreement near numerical precision when the Jacobian
is present and large errors when it is removed (Section 4.1).

\subsubsection{3.2 Exact local Gaussian
packet}\label{exact-local-gaussian-packet}

Consider a strict spacetime Gaussian in a nonsingular local affine
camera-ray chart \(s=(r,z)\), with \(r=(u,v,\tau)\), mean
\((\mu_r,\mu_z)\), and covariance blocks
\(\Sigma_{rr},\Sigma_{rz},\Sigma_{zr},\Sigma_{zz}\).

Gaussian marginalization gives the UVT footprint

\[
  r\sim\mathcal N(\mu_r,\Sigma_{rr}),
  \qquad Q_r=\Sigma_{rr}^{-1},
  \tag{3}
\]

while conditioning retains the depth information needed for visibility:

\[
\begin{aligned}
  \widehat z(r)
    &=\mu_z+\Sigma_{zr}Q_r(r-\mu_r),\\
  \nu_z
    &=\Sigma_{zz}-\Sigma_{zr}Q_r\Sigma_{rz}>0.
\end{aligned}
\tag{4}
\]

The strict inequality is the Schur complement of a positive-definite
covariance. Thus one source object yields an anisotropic sensor-time
footprint, an affine conditional-depth plane, and a positive
conditional-depth variance. World Tubes marginalizes depth for fast
footprint evaluation but does not discard the conditional packet used to
certify order or invoke a retained-depth fallback.

Equation (3) is exact only for an affine chart with the declared fiber
measure and integration domain. A moving pinhole camera is nonlinear
after perspective division, so long windows require either certified
local approximation or the projective construction below. Complete
covariance and equivalent precision derivations belong in the
supplement.

\subsubsection{3.3 Projective trace
domains}\label{projective-trace-domains}

For homogeneous primitive center \(X_i^h(\tau)\) and camera matrix
\(P(\tau)\), retain the undivided trace and divide only after
certification:

\[
\begin{aligned}
  h_i(\tau)&=P(\tau)X_i^h(\tau)
  =\bigl(h_{u,i}(\tau),h_{v,i}(\tau),h_{z,i}(\tau)\bigr),\\
  (u_i,v_i,d_i)&=\left(\frac{h_{u,i}}{h_{z,i}},
          \frac{h_{v,i}}{h_{z,i}},h_{z,i}\right).
\end{aligned}
  \tag{5}
\]

On a bounded interval \(I_\ell\), normalized time \(s\in[-1,1]\)
parameterizes the supported degree-one or degree-two trace family

\[
  h_{k,i}(s)=\sum_{q=0}^{p}a_{k,i,q}s^q,
  \qquad k\in\{u,v,z\}.
  \tag{6}
\]

The compiler accepts an interval only when four contracts hold:

\begin{enumerate}
\def\labelenumi{\arabic{enumi}.}
\tightlist
\item
  \textbf{Projective validity:} \(h_{z,i}\) retains the required
  physical sign and remains separated from zero over the complete
  interval.
\item
  \textbf{Trace accuracy:} divided UV and conditional depth remain
  within their declared residual tolerances.
\item
  \textbf{Support validity:} the footprint plus approximation padding
  conservatively fixes tile membership and active-time gates.
\item
  \textbf{Visibility validity:} every support-overlapping pair has a
  certified order, an accepted commutation residual, or an explicit
  fallback label.
\end{enumerate}

The implemented denominator certificate is continuous for its quadratic
family. General UV and depth residuals are bounded on the declared probe
set; we do not promote that sampled test into an unrestricted
continuous-time theorem. A failed interval is split at a certified event
when available and at a deterministic midpoint otherwise. A minimum-size
interval that still fails is routed to a reason-labelled fallback.

\begin{figure}
\centering
\includegraphics[width=1\linewidth,height=\textheight,keepaspectratio]{/Users/nicholasbardy/git/gsplats_browser/dynaworld/tmp/pdfs/world_tubes_projective_compiler.png}
\caption{A projective trace is divided only after denominator, residual,
support, and visibility certificates all pass. Failure splits the
bounded interval or routes it to an explicit reference fallback; the
compiler never extrapolates through an unresolved pole or order
event.}\label{fig:world-tubes-projective-compiler}
\end{figure}

\subsubsection{3.4 Visibility-stratified trace
atlas}\label{visibility-stratified-trace-atlas}

The compiler emits cells

\[
  \mathcal A_\Gamma=
  \{(C_\ell,S_\ell,\Phi_\ell,\Pi_\ell,E_\ell)\}_{\ell=1}^{L},
  \tag{7}
\]

where \(C_\ell\subset B\) is a certified sensor-time domain, \(S_\ell\)
its active primitive set, \(\Phi_\ell\) the trace records, \(\Pi_\ell\)
the compiled order or partial-order certificate, and \(E_\ell\) the
support, error, and fallback metadata.

For primitive \(i\), let \([z_i^-(y),z_i^+(y)]\) conservatively bound
conditional depth. A pair has a fixed order when

\[
  z_i^+(y)<z_j^-(y)
  \quad\text{or}\quad
  z_j^+(y)<z_i^-(y)
  \qquad\forall y\in C_\ell.
  \tag{8}
\]

Otherwise the cell is split, the pair is shown to commute within
tolerance, or the cell falls back. Swapping two alpha-composited
contributors changes their local color by

\[
  \delta_{ij}(y)
  =\alpha_i(y)\alpha_j(y)\bigl(c_i(y)-c_j(y)\bigr).
  \tag{9}
\]

This gives a baseline-relative error certificate for accepted adjacent
swaps. It does not claim that representative-depth alpha compositing is
physically exact radiative transfer. It states when the chosen frozen
renderer can reuse one visibility topology.

The final color inside a cell is evaluated in compiled front-to-back
order:

\[
  I(y)=\sum_m
  \left[\prod_{n<m}(1-\alpha_{\pi_n}(y))\right]
  \alpha_{\pi_m}(y)c_{\pi_m}(y).
  \tag{10}
\]

Finite exposure integrates Eq. (10) over shutter time; rolling shutter
changes the sampling map from row position to time. Neither requires
recompiling the world while the requested samples remain in the same
certified cells.

\subsubsection{3.5 Compiled adjoint}\label{compiled-adjoint}

Let \(\theta\) denote the camera-independent world parameters, let
\(\kappa\) collect discrete support, bin, event, order, and fallback
decisions, and let

\[
  \phi=C_\Gamma(\theta;\kappa),\qquad
  I=R(\phi,\kappa).
\]

Within a fixed-topology stratum, the world gradient is

\[
  \nabla_\theta\mathcal L
  =D_\theta C_\Gamma(\theta;\kappa)^\top
   D_\phi R(\phi,\kappa)^\top
   \nabla_I\ell(I,I^\star).
  \tag{11}
\]

The first reverse map is the interval-renderer VJP; the second
accumulates through projective coefficients, gauge transformations, and
the local Gaussian packet into the same world parameters used by replay.
Event-boundary derivatives are not hidden inside Eq. (11). When a
parameter update changes \(\kappa\), the compiler must refresh the
affected structure or use a separately declared boundary estimator.

\subsubsection{3.6 Work boundary}\label{work-boundary}

For \(F\) requested times and materialized interaction count \(K\), a
coarse decomposition is

\[
\begin{aligned}
W_{\mathrm{replay}}&=W_{\mathrm{shade}}(K)+\sum_{f=1}^{F}W_{\mathrm{world}}^f,\\
W_{\mathrm{tube}}&=W_{\mathrm{shade}}(K)+W_{\mathrm{compile}}+W_{\mathrm{eval}}(F).
\end{aligned}
\tag{12}
\]

Both routes produce \(FHW\) samples and pay for real shading
interactions. World Tubes targets the repeated projection, conservative
support, tile membership, visibility metadata, and corresponding
backward tape. For a fixed bounded camera program whose chart and event
structure remains stable as time sampling is densified, the persistent
atlas state can remain fixed while replay metadata grows with \(F\).
Near poles, support churn, dense order crossings, or fallback-heavy
regions can instead make the compiled structure grow. We claim
conditional camera-program amortization, not a universal sublinear bound
on end-to-end rendering.

\subsection{4. Implementation}\label{implementation}

Our implementation retains the STAR UVT/projective-interval execution
lineage but changes the compiler semantics. The compiler stores
homogeneous trace coefficients, optional temporal-opacity and
spatial-precision terms, conditional-depth information, active-time
intervals, tile membership, and certified cell order. Accepted cells are
packed once into tile-time bins and evaluated by an interval Metal
forward pass. A custom autograd boundary applies the direct interval VJP
and then the differentiable compiler map of Eq. (11).

Support, bin membership, interval splits, order decisions, and fallback
labels are discrete during one adjoint block. Trace coefficients,
Gaussian source parameters, opacity, precision, temporal terms, and
appearance remain differentiable. The implemented moving-camera source
is bounded to one first-order or projective chart at a time; complete
repeated-orbit chart transitions are outside the claim. The
retained-depth branch currently covers the scoped affine Gaussian
fallback rather than nonlinear projective fibers.

The supplement records source parameterizations, amplitude conventions,
Beer--Lambert support and VJPs, compiler tolerances, kernel entry
points, and the exact fixed-topology refresh policy. These details are
omitted here because they do not change the causal comparison: replay
and compilation start from the same world and differ only in whether
world-to-trace work is repeated per frame.

\subsection{5. Experiments}\label{experiments}

\subsubsection{5.1 Certified bounded
correctness}\label{certified-bounded-correctness}

The certified suite derives its rows from byte-pinned source artifacts
without loading the training runtime. It tests gauge changes, replay
equivalence, visibility crossings, exposure and rolling-shutter
evaluation, mixed fallback, and fixed-topology gradients. Table 1
reports the accepted bounded scope; no timing row is inferred from these
fixtures.

The gauge ablation is especially diagnostic. Across five sensor-time
points, ordinary- and log-depth integration agree to at most
\(3.50\times10^{-13}\) relative error with the Jacobian; removing it
gives at least \(0.600\) relative error. Gradients agree to
\(2.33\times10^{-12}\) with the Jacobian, while the no-Jacobian variants
err by at least \(0.592\). The visibility ablation isolates a different
failure: one unstratified interval produces \(0.186742\) error at an
order crossing, whereas event stratification reproduces the live-order
oracle exactly on the fixture.

\subsubsection{5.2 Same-representation frame
scaling}\label{same-representation-frame-scaling}

The structural fixture evaluates \(F\in\{4,8,16,32,64,128\}\). Its fixed
trace tensor volume remains constant while the per-frame replay
construction grows by \(32\times\) from the first to the last row; the
final trace-count ratio is \(0.03125\). This is evidence of structural
reuse, not a runtime, storage, or peak-memory claim: logical tensors
exclude packed topology, allocator overhead, transient buffers, and
unavoidable pixel interactions.

The causal learned-scene experiment freezes one checkpoint
\(\theta^\star\), camera program, target samples, precision, background,
shading, loss, and pixel batch. Replay lowers that world independently
for every selected time; World Tubes compiles the same world once for
the complete bounded program. Both routes evaluate identical targets
over the same physical interval and return gradients to the same world
parameters.

\subsubsection{5.3 Public evaluation
contract}\label{public-evaluation-contract}

Neural 3D Video supplies the positive calibrated multiview setting
(\citeproc{ref-li2022_neural3dvideo}{Li et al. 2022}); D-NeRF is a
labelled posed-frame control (\citeproc{ref-pumarola2021_dnerf}{Pumarola
et al. 2021}), not evidence for multi-camera amortization. The primary
public compiler comparison is the frozen same-world protocol above. A
separate matched-schedule matrix compares World Tubes, WorldFoam, and
dynamic 3DGS as representation-and-cost context, not compiler causality.
Its evaluator fixes inputs, split, background, and metric aggregation,
and reports storage, rasterized samples, peak memory, and wall time
alongside image quality.

\subsubsection{5.4 Central ablations}\label{central-ablations}

The bounded evidence separates three mechanisms from a generic ``faster
STAR'' result. Retaining the gauge Jacobian changes value and gradient
invariance by more than eleven orders of magnitude; event stratification
removes the observed depth-order-crossing error; and fixed trace state
reduces the \(F=128\) trace count to \(3.125\%\) of replay without
claiming that output shading disappears. Clean closure and learned-scene
artifacts add the affine/projective, runtime, backward, peak-memory, and
break-even ablations. Ordered transfer counts as a positive result only
where its certificate is selective; all-retained routing remains a
negative control.

\subsection{6. Discussion and
limitations}\label{discussion-and-limitations}

World Tubes targets known or low-dimensional camera programs with
repeated temporal samples. Random views offer no amortization;
visibility chaos, near-camera splats, wide fields of view, or support
churn may erase it.

The claim is local: the exact Schur complement requires a regular affine
chart, projective traces cover bounded certified segments rather than
one complete \(360^\circ/720^\circ\) chart, and the adjoint is exact
inside a fixed discrete topology. Event boundaries require refresh or a
declared estimator. Conditional Gaussian depth does not represent
arbitrary interleaved colored density; ambiguous cells may use ordered
transfer, while general cellular fields and adaptive retained-depth
quadrature remain outside the central claim.

\subsection{7. Conclusion}\label{conclusion}

World Tubes compiles dynamic Gaussian scenes into gauge-correct,
event-certified traces whose fixed-topology adjoint shares world-side
work across time. Certified fixtures establish invariance, failure
detection, gradient parity, and structural reuse.

\subsection{Reproducibility statement}\label{reproducibility-statement}

The repository records checked-in protocols, deterministic seed lists,
dataset and split contracts, source and native-binary hashes,
isolated-run summaries, raw timing samples, retained artifacts, and
Torch-free evidence verifiers. The causal comparison binds replay and
compiled routes to one checkpoint and one camera program. The supplement
will enumerate compiler tolerances, parameter conventions, evaluation
formulas, hardware, warmup and repeat policy, and the commands that
regenerate every accepted table and figure.

\subsection{References}\label{references}

\protect\phantomsection\label{refs}
\begin{CSLReferences}{1}{1}
\bibitem[\citeproctext]{ref-hou2025_sortfree}
Hou, Qiqi, Randall Rauwendaal, Zifeng Li, et al. 2025. {``{Sort-Free
Gaussian Splatting via Weighted Sum Rendering}.''} \emph{International
Conference on Learning Representations}.
\url{https://openreview.net/forum?id=y8uPsxR8PN}.

\bibitem[\citeproctext]{ref-kerbl2023}
Kerbl, Bernhard, Georgios Kopanas, Thomas Leimkühler, and George
Drettakis. 2023. {``{3D Gaussian Splatting for Real-Time Radiance Field
Rendering}.''} \emph{ACM Transactions on Graphics} 42 (4): 1--14.
\url{https://doi.org/10.1145/3592433}.

\bibitem[\citeproctext]{ref-koo2026_blending}
Koo, Junseo, Jinseo Jeong, and Gunhee Kim. 2026. {``{Gaussian Blending:
Rethinking Alpha Blending in 3D Gaussian Splatting}.''}
\emph{Proceedings of the AAAI Conference on Artificial Intelligence} 40:
5744--54. \url{https://doi.org/10.1609/aaai.v40i7.37495}.

\bibitem[\citeproctext]{ref-li2022_neural3dvideo}
Li, Tianye, Mira Slavcheva, Michael Zollhöfer, et al. 2022. {``{Neural
3D Video Synthesis From Multi-View Video}.''} \emph{Proceedings of the
IEEE/CVF Conference on Computer Vision and Pattern Recognition},
5521--31.

\bibitem[\citeproctext]{ref-li2024_spacetime}
Li, Zhan, Zhang Chen, Zhong Li, and Yi Xu. 2024. {``{Spacetime Gaussian
Feature Splatting for Real-Time Dynamic View Synthesis}.''}
\emph{Proceedings of the IEEE/CVF Conference on Computer Vision and
Pattern Recognition}, 8508--20.

\bibitem[\citeproctext]{ref-luiten2024}
Luiten, Jonathon, Georgios Kopanas, Bastian Leibe, and Deva Ramanan.
2024. {``{Dynamic 3D Gaussians: Tracking by Persistent Dynamic View
Synthesis}.''} \emph{2024 International Conference on 3D Vision (3DV)},
800--809. \url{https://doi.org/10.1109/3DV62453.2024.00044}.

\bibitem[\citeproctext]{ref-mcguire2013_weighted_oit}
McGuire, Morgan, and Louis Bavoil. 2013. {``{Weighted Blended
Order-Independent Transparency}.''} \emph{Journal of Computer Graphics
Techniques} 2 (2): 122--41.
\url{https://www.jcgt.org/published/0002/02/09/}.

\bibitem[\citeproctext]{ref-pumarola2021_dnerf}
Pumarola, Albert, Enric Corona, Gerard Pons-Moll, and Francesc
Moreno-Noguer. 2021. {``{D-NeRF: Neural Radiance Fields for Dynamic
Scenes}.''} \emph{Proceedings of the IEEE/CVF Conference on Computer
Vision and Pattern Recognition}, 10318--27.

\bibitem[\citeproctext]{ref-radl2024_stopthepop}
Radl, Lukas, Michael Steiner, Mathias Parger, Alexander Weinrauch,
Bernhard Kerbl, and Markus Steinberger. 2024. {``{StopThePop: Sorted
Gaussian Splatting for View-Consistent Real-time Rendering}.''}
\emph{ACM Transactions on Graphics} 43 (4): 1--17.
\url{https://doi.org/10.1145/3658187}.

\bibitem[\citeproctext]{ref-seiskari2024}
Seiskari, Otto, Jerry Ylilammi, Valtteri Kaatrasalo, et al. 2025.
{``{Gaussian Splatting on the Move: Blur and Rolling Shutter
Compensation for Natural Camera Motion}.''} \emph{Computer Vision --
ECCV 2024}, Lecture notes in computer science, vol. 15129: 160--77.
\url{https://doi.org/10.1007/978-3-031-73209-6_10}.

\bibitem[\citeproctext]{ref-wu2024_4dgs}
Wu, Guanjun, Taoran Yi, Jiemin Fang, et al. 2024. {``{4D Gaussian
Splatting for Real-Time Dynamic Scene Rendering}.''} \emph{Proceedings
of the IEEE/CVF Conference on Computer Vision and Pattern Recognition},
20310--20.

\bibitem[\citeproctext]{ref-wu2025_3dgut}
Wu, Qi, Janick Martinez Esturo, Ashkan Mirzaei, Nicolas Moënne-Loccoz,
and Zan Gojcic. 2025. {``{3DGUT: Enabling Distorted Cameras and
Secondary Rays in Gaussian Splatting}.''} \emph{Proceedings of the
IEEE/CVF Conference on Computer Vision and Pattern Recognition},
26036--46.

\bibitem[\citeproctext]{ref-yang2024_native}
Yang, Zeyu, Zijie Pan, Xiatian Zhu, et al. 2024. {``{4D Gaussian
Splatting: Modeling Dynamic Scenes with Native 4D Primitives}.''}
\emph{arXiv Preprint arXiv:2412.20720}.
\url{https://arxiv.org/abs/2412.20720}.

\bibitem[\citeproctext]{ref-yang2024_deformable}
Yang, Ziyi, Xinyu Gao, Wen Zhou, Shaohui Jiao, Yuqing Zhang, and
Xiaogang Jin. 2024. {``{Deformable 3D Gaussians for High-Fidelity
Monocular Dynamic Scene Reconstruction}.''} \emph{Proceedings of the
IEEE/CVF Conference on Computer Vision and Pattern Recognition},
20331--41.

\bibitem[\citeproctext]{ref-yu2024_mipsplatting}
Yu, Zehao, Anpei Chen, Binbin Huang, Torsten Sattler, and Andreas
Geiger. 2024. {``{Mip-Splatting: Alias-free 3D Gaussian Splatting}.''}
\emph{Proceedings of the IEEE/CVF Conference on Computer Vision and
Pattern Recognition}, 19447--56.

\bibitem[\citeproctext]{ref-zwicker2001_ewa_volume}
Zwicker, Matthias, Hanspeter Pfister, Jeroen van Baar, and Markus Gross.
2001a. {``{EWA Volume Splatting}.''} \emph{Proceedings of IEEE
Visualization}, 29--36.
\url{https://www.merl.com/publications/TR2001-27}.

\bibitem[\citeproctext]{ref-zwicker2001_surface}
Zwicker, Matthias, Hanspeter Pfister, Jeroen van Baar, and Markus Gross.
2001b. {``{Surface Splatting}.''} \emph{Proceedings of the 28th Annual
Conference on Computer Graphics and Interactive Techniques}, 371--78.
\url{https://doi.org/10.1145/383259.383300}.

\end{CSLReferences}

\appendix

\subsection{A. Ordered ray transfer at the noncommutation
boundary}\label{a.-ordered-ray-transfer-at-the-noncommutation-boundary}

One scalar depth order is inadequate when thick, differently colored
density profiles interleave along a ray. The scoped retained-depth
extension represents this case by the ordered optical-transfer product

\[
  M_y=\mathcal P\exp\!\left(\int A_y(z)\,\mathrm dz\right),\qquad
  A_y(z)=
  \begin{bmatrix}
    -\sigma_y(z)I_3 & \sigma_y(z)c_y(z)\\
    0 & 0
  \end{bmatrix}.
  \tag{13}
\]

The product integral is invariant under an orientation-preserving depth
reparameterization when its physical-length Jacobian is included. Here
it is a bounded fallback for Gaussian depth fibers, not a replacement
for the fast certified compositor and not a claim about a general
cellular field. The full supplement must include the transfer
derivation, VJP, certificate, selective fixture, and dense-scene
negative control.

\subsection{Supplementary-material
roadmap}\label{supplementary-material-roadmap}

The supplement should preserve, rather than delete, the details removed
from the nine-page main-paper budget:

\begin{enumerate}
\def\labelenumi{\arabic{enumi}.}
\tightlist
\item
  notation, camera-ray bundle construction, and gauge-invariance proof;
\item
  full SPD(4) covariance and equivalent precision derivations;
\item
  amplitude conventions and Beer--Lambert support/VJP derivations;
\item
  projective trace coefficients, continuous denominator certificate, and
  compiler pseudocode;
\item
  visibility commutation bound and ordered ray-transfer derivation;
\item
  fixed-topology adjoint proof and complete work accounting;
\item
  source parameterizations, kernel interfaces, topology-refresh policy,
  and numerical tolerances;
\item
  dataset contracts, splits, schedules, timing and memory methodology;
\item
  full closure/death diagnostics, all public rows, additional scenes and
  seeds, qualitative examples, and failure-bearing negative controls.
\end{enumerate}

\end{document}
